\section{Identidades de transferencia, memoria y geometría}
\label{app:wolfram}

Las realizaciones de Lur'e permiten expresar el campo mediante matrices y
no linealidades explícitas. Su consistencia se comprueba sustituyendo la
realización en el campo original y simplificando la diferencia.

\subsection{Identidades verificadas}

Las identidades fundamentales son:
\begin{enumerate}
 \item igualdad entre el campo original y la forma de Lur'e
 \(PX+b\psi(r^{\T}X)\);
 \item identidad de la transferencia
 \(r^{\T}(\zeta I-P)^{-1}b=-\alpha N(\zeta)/\Delta_\ell(\zeta)\);
 \item identidad del determinante del sistema auxiliar obtenida mediante el
 lema del determinante;
 \item ecuación escalar de los equilibrios, jacobiano y polinomio
 característico;
 \item forma cerrada de la función descriptiva;
 \item sustitución de \(\zeta=(\ii\omega)^q\) en la rama principal;
 \item residuos numéricos de los equilibrios y márgenes de Matignon.
\end{enumerate}

En la familia con saturación también se comprueba la transformación modal
\begin{equation}
P_0S=SH_q,\qquad
H_q=
\begin{bmatrix}
\zeta_r&-\zeta_i&0\\
\zeta_i& \zeta_r&0\\
0&0&-d
\end{bmatrix},
\label{eq:wolfram_similitud}
\end{equation}
y se calcula la semilla como \(A_0S_{\cdot1}\). En la familia con arctangente
se integra simbólicamente la primera armónica y se simplifica hasta la forma
cerrada de \cref{eq:df_arctan}. Para c590, la semilla se obtiene mediante búsqueda entera acotada y
refinamiento de Caputo. La identidad de función descriptiva caracteriza la
familia arctangente, independientemente del procedimiento de selección.

\subsection{Composición temporal y memoria}

La expansión
\[
 E_q(\lambda h^q)=1+\frac{\lambda h^q}{\Gamma(1+q)}+O(h^{2q})
\]
permite comparar dos evoluciones de duración \(h\). Para \(0<q<1\) y
\(\lambda\ne0\),
\[
 E_q\bigl(\lambda(2h)^q\bigr)-E_q(\lambda h^q)^2
 =\frac{\lambda(2^q-2)}{\Gamma(1+q)}h^q+O(h^{2q}).
\]
El coeficiente principal es distinto de cero, pues \(2^q<2\). Por tanto,
la diferencia no se anula para todo \(h>0\) suficientemente pequeño y la ley
de semigrupo puntual falla. Una evaluación concreta ilustra su magnitud.

Para
\[
 q=\frac45,\qquad \lambda=-\frac45,\qquad
 t=\frac35,\qquad s=\frac7{10},
\]
se obtuvo
\begin{align}
 E_q\!\left(\lambda(t+s)^q\right)
 &=
 0.3911882069642177396\ldots,\\
 E_q(\lambda t^q)E_q(\lambda s^q)
 &=
 0.3209479300285948495\ldots,
\end{align}
con residuo \(0.0702402769356228901\ldots\). Por tanto, la solución de Mittag--Leffler no satisface la ley de semigrupo
temporal sobre el estado puntual. Para \(q=1\), la identidad exponencial
restablece dicha ley. En cambio, la integral fraccionaria satisface
\[
 I^\beta I^\alpha t^\mu=I^{\alpha+\beta}t^\mu
\]
para \(\alpha,\beta>0\) y \(\mu>-1\), con terminal inferior cero. En efecto,
el cambio \(s=tu\) y la integral beta dan
\begin{align*}
 I^\alpha t^\mu
 &=\frac{t^{\mu+\alpha}}{\Gamma(\alpha)}
   \int_0^1(1-u)^{\alpha-1}u^\mu\,\mathrm du
 =\frac{\Gamma(\mu+1)}{\Gamma(\mu+\alpha+1)}t^{\mu+\alpha},\\
 I^\beta I^\alpha t^\mu
 &=\frac{\Gamma(\mu+1)}{\Gamma(\mu+\alpha+1)}
   \frac{\Gamma(\mu+\alpha+1)}{\Gamma(\mu+\alpha+\beta+1)}
   t^{\mu+\alpha+\beta}
 =I^{\alpha+\beta}t^\mu.
\end{align*}
Esta composición se refiere al
\emph{orden del operador}; no expresa evolución temporal.

En el campo escalar \(f(x)=\lambda x\), las dos primeras iteraciones de
Picard--Caputo producen
\(x_0[1+\lambda h^q/\Gamma(1+q)+\lambda^2h^{2q}/\Gamma(1+2q)]\).
Al restarlas de la solución \(x_0E_q(\lambda h^q)\) y dividir por
\(h^{3q}\), el límite es \(x_0\lambda^3/\Gamma(1+3q)\), directamente de
la serie. El refinamiento \(h=1/10,1/20,1/40,1/80\) reproduce ese límite.

En geometría entera, un equilibrio no degenerado \(e\) tiene índice
\(\operatorname{sgn}\det J_f(e)\). Los valores de los casos considerados son:
\begin{center}
\begin{tabular}{@{}lll@{}}
\toprule
Caso & determinantes evaluados & índices locales \\
\midrule
Chua PWL & \(-84.035507517056\), \(15.037737580544\)
 & \((-1,1,1)\)\\
MAVPD & \(2000,-4000,-4000\) & \((1,-1,-1)\)\\
PLL & \(2758.72946512,-2758.72946512\) & \((1,-1)\)\\
\bottomrule
\end{tabular}
\end{center}

Estas comprobaciones determinan identidades, residuos e índices locales. Un
certificado global de grado de Brouwer requeriría una región acotada y la
ausencia de ceros en su frontera. Las trayectorias, cuencas, atractores,
ocultedad, entropía y caos requieren las campañas dinámicas y las hipótesis de
atracción y perfiles admisibles de la \cref{sec:geometria-topologia-caputo}.
En particular, la \cref{prop:geofo-full-space-obstruction} descarta la
atracción de una vecindad abierta en todo el espacio de términos libres; una
identidad simbólica no elimina esa obstrucción.

\subsection{Propiedades del localizador geométrico--topológico}

La expansión de Picard--Caputo y las leyes de transformación de la
\cref{sec:localizacion-geometrico-topologica} conducen a las siguientes
propiedades:
\begin{enumerate}
 \item la covariancia de la curva de arranque transportada y, por separado,
 el defecto al recalcular sus coeficientes a partir del campo transformado:
 $(a_q-b_q)D^2g[F,F]$. El contraejemplo con $q=1/2$ produce
 $(0,4-8/\pi)\ne0$, mientras la aceleración covariante de la curva transportada
 es nula;
 \item las identidades racionales de $\nabla V\cdot F$ y del resto no negativo
 en la \cref{prop:geofo-lorenz-boundedness}. La comparación de Caputo y la
 continuación global se justifican en la demostración de esa proposición;
 \item las superficies regionales de aceleración de Chua PWL;
 \item las superficies de aceleración de MAVPD y la inclusión de sus puntos
 perpetuos en la curva conectante;
 \item la aceleración en forma compañera y los seis menores de rango de
 Kalman--Fitts;
 \item las superficies de aceleración y la ecuación de curva conectante del
 PLL;
 \item la identidad del segundo invariante KCC del PLL y los signos
 \[
 P(\theta_f,0)=-1642.024105242218\ldots<0,
 \qquad
 P(\theta_s,0)=3069.171079266960\ldots>0;
 \]
 \item los determinantes, simetrías, raíces y residuos de Chua arctangente,
 Lorenz generalizado, Rabinovich--Fabrikant y Muñoz--Pacheco.
\end{enumerate}
El signo de una curvatura KCC, la anulación de un menor o una raíz de
\(J_FF\) son propiedades algebraicas locales. Por separado no demuestran que
exista un atractor, que sea caótico ni que su cuenca sea oculta. Esas preguntas
corresponden a la campaña dinámica de la \cref{sec:campana-gt}.

La validación algebraica comprueba las operaciones que enlazan
\cref{eq:chua_no_suave,eq:chua_arctan} con las ecuaciones auxiliares. La
caracterización dinámica se completa con los PVI, los controles de paso y los
sondeos descritos en \cref{sec:integracion_continuacion,sec:protocolos,sec:resultados}.
